\scanpage{167}\mspage{112}
pleine fidélité de $i$, sont valables dans le cas $\deleted{\unsure{B}}A$, quand on renforce l'hypothèse A (relative à un $S$ particulier) en une hypothèse plus restrictive et plus compréhensive, que voici.

A$'$) Le foncteur $C:\unsure{M|\mathcal M}_{/S}\longrightarrow M_{/S}$ et \unsure{le|les} \deleted{\unsure{immersions}} \unsure{morphisme|morphismes} \inserted{\unsure{fonctoriel|fonctoriels}} $i_X:X\longrightarrow C(X)$ sont construits ``à la \unsure{Cartan-Serre}'' à partir d'un foncteur défini sur $M$ \deleted{tout} entier
\[
M\xrightarrow{\mathrm{Ch}}M,
\]
du type ``chemins'', i.e. muni en plus des données suivantes
\[
X\xrightarrow{\alpha_X}\mathrm{Ch}(X)\xrightarrow{p_X=(\sigma_X,\beta_X)}X\times X\tag*{a)}
\]
(structure de source et but des chemins, et de chemins \deleted{\unsure{identiques}}\inserted{triviaux}), et

b) Une composition des chemins, sans exiger ici l'associativité de la composition, mais seulement que les chemins triviaux soient \unsure{neutres|unités} bilatères pour la composition.

De plus, on suppose que pour tout $X$
\[
\left\{\begin{aligned}
&\alpha_X\in W,\quad p_X\in\mathrm{Fib}\\
&(X\longrightarrow e)\in\mathrm{Fib}.
\end{aligned}\right.\tag*{\deleted{\unsure{(3)}}}
\]
% archive sheet number: 166
